% !TEX root = ../../hw07.tex
% Source: Lucas's supplied typed solutions, October 30, 2026, 1.1--1.2.
% Correct the linearity-test typo (u,v) to (u,w), as in his calculation.
\begin{proof}[Solution]
  \textbf{1.1.} Testing positive definiteness gives
  \[
    (v,v)=v^1v^2+v^2v^1=2v^1v^2.
  \]
  The real coordinates \(v^1,v^2\) can have opposite signs, so this
  expression can be negative even when \(v\ne0\). Thus the form is
  not positive definite and is not an inner product.

  \textbf{1.2.} First, the form is symmetric because
  \[
    \begin{aligned}
      (w,v)&=2w^1v^1-w^1v^2-w^2v^1+2w^2v^2\\
        &=(v,w).
    \end{aligned}
  \]
  To test positive definiteness, calculate
  \[
    \begin{aligned}
      (v,v)&=2(v^1)^2-2v^1v^2+2(v^2)^2\\
        &=\bigl((v^1)^2-2v^1v^2+(v^2)^2\bigr)
          +(v^1)^2+(v^2)^2\\
        &=(v^1-v^2)^2+(v^1)^2+(v^2)^2.
    \end{aligned}
  \]
  A sum of real squares is nonnegative and vanishes only when every
  squared term is zero. Hence \((v,v)=0\) exactly when \(v=0\), and
  \((v,v)>0\) whenever \(v\ne0\).

  For linearity in the first argument, let \(a\in\R\) and
  \(u,v,w\in\R^2\) be arbitrary. Then
  \[
    \begin{aligned}
      (av+u,w)
        &=2(av^1+u^1)w^1-(av^1+u^1)w^2\\
        &\quad -(av^2+u^2)w^1+2(av^2+u^2)w^2\\
        &=a\bigl(2v^1w^1-v^1w^2-v^2w^1+2v^2w^2\bigr)\\
        &\quad +\bigl(2u^1w^1-u^1w^2-u^2w^1+2u^2w^2\bigr)\\
        &=a(v,w)+(u,w).
    \end{aligned}
  \]
  By symmetry,
  \[
    (v,aw+u)=(aw+u,v)=a(w,v)+(u,v)=a(v,w)+(v,u).
  \]
  Thus the form is bilinear, symmetric, and positive definite, so it
  is an inner product.
\end{proof}
